\section{Einleitung}\label{sec:einleitung}

\subsection{Ausgangslage}\label{sec:ausgangslage}

Jass ist ein verbreitetes Kartenspiel in der Schweiz. Am Ende einer Runde
werden die gewonnenen Stiche einer Partei gezählt, indem die Karten
nacheinander abgelegt und ihre Punktwerte addiert werden. Das kostet Zeit
und ist besonders für ungeübte Spielende fehleranfällig.

Smartphones haben heute leistungsfähige Kameras und Recheneinheiten für
maschinelles Lernen. Damit lässt sich ein Kartenstapel direkt am Spieltisch
filmen und auswerten, ohne einen Server zu benötigen. In dieser Arbeit wird
ein Modell trainiert, das in einem Videostream jeweils die oberste Karte
eines Stapels erkennt. Die App soll daraus die Punkte der Runde berechnen.

Ursprünglich war vorgesehen, die Trainingsbilder selbst zu fotografieren
und von Hand zu labeln. Für die benötigte Menge erwies sich das früh als
unrealistisch. Deshalb werden die Trainingsbilder synthetisch aus Fotos der
einzelnen Karten erzeugt. Diese Kurskorrektur ist ein wichtiger Teil der
Arbeit und wird in Abschnitt~\ref{sec:verworfen} reflektiert.

\subsection{Problemstellung}\label{sec:problemstellung}

Erkannt werden soll nur die oberste Karte eines Stapels überlappender
Karten. Darunterliegende, teilweise sichtbare Karten dürfen nicht gezählt
werden. Die Erkennung muss zuverlässig und schnell genug sein, damit die
App in einer realen Spielsituation genutzt werden kann. Wer die Karten
zählt, soll das Ergebnis nicht nach jeder Erkennung selbst kontrollieren
müssen.

Die Aufnahmebedingungen variieren: Die Kamera wird von Hand geführt, der
Abstand und der Winkel ändern sich, und das Licht kann von einer
Deckenlampe, vom Tageslicht oder von der Taschenlampe des Smartphones
kommen. Auch die Unterlage wechselt zwischen Jassteppich, Holz und Stein.
In vielen Frames ist keine eindeutig lesbare oberste Karte zu sehen, etwa
wenn eine Karte geworfen wird, verwischt oder den Stapel verdeckt.

Für solche Situationen gilt die Prämisse dieser Arbeit: \emph{Lieber
nichts als falsch.} Eine übersehene Karte kann erneut vor die Kamera
gehalten werden. Eine falsch gezählte Karte bleibt dagegen möglicherweise
unbemerkt und verfälscht das Resultat. Gerade weil ein Videostream mehrere
Frames derselben Karte liefert, kann die App kurze Lücken überbrücken.

Der Engpass liegt bei den Trainingsdaten. Ultralytics nennt 1500 Bilder je
Klasse als Richtwert für eigene Datensätze \cite{ultralytics2025tips}.
Schon für die 36 Karten eines Blatts wären das 54'000 Bilder. Diese Menge
von Kartenstapeln lässt sich im Rahmen der Transferarbeit weder sinnvoll
fotografieren noch von Hand labeln.

\subsection{Ziele und Forschungsfragen}\label{sec:ziele}

Ziel ist ein mobiler Prototyp, der die oberste Karte mithilfe eines auf
synthetischen Bildern trainierten Modells erkennt und die Jasspunkte eines
Stapels zählt. Daraus ergeben sich zwei Forschungsfragen:

\begin{description}
  \item[FF1] Welcher Erkennungsansatz (einstufiger Detektor, zweistufige
    Erkennung mit orientierter Box und Klassifikation oder reine
    Bildklassifikation) bietet auf dem Smartphone die beste Balance
    zwischen Genauigkeit und Rechenaufwand?
  \item[FF2] Wie gut übertragen sich Modelle, die ausschliesslich auf
    synthetisch erzeugten Bildern trainiert wurden, auf reale Fotos?
\end{description}

Die ursprüngliche Themeneingabe sah auch einen Vergleich zwischen realen
und synthetischen Trainingsbildern vor. Da keine ausreichend grosse Menge
realer Trainingsbilder erstellt werden konnte, wird stattdessen der
Transfer von synthetischen Trainingsbildern auf reale Aufnahmen geprüft.
Ebenso entfiel der vorgesehene Vergleich verschiedener Modellgrössen.
Alle Modelle verwenden die kleinste Grösse n, weil sie auf dem Smartphone
in Echtzeit läuft. An die Stelle dieses Vergleichs trat der Vergleich der
drei Erkennungsansätze. Geschwindigkeit und Modellgrösse wurden für den
gewählten Ansatz auf vier Geräten gemessen
(Abschnitt~\ref{sec:app}).

\subsection{Abgrenzung}\label{sec:abgrenzung}

Die App erkennt die oberste Karte und zählt den Stapel einer Partei, in der
Praxis häufig den kleineren der beiden Stapel. Aus der Gesamtpunktzahl
der Karten und der Angabe zum letzten Stich berechnet sie auch die Punkte
der anderen Partei. Ein vollständiges Scoring mit Weis und Stöck ist
nicht Teil der Arbeit.

Die Karten liegen flach auf dem Tisch. Gekippte oder stehende Karten sowie
Hände im Bild werden im Generator nicht modelliert. Unterstützt werden
zwei Blätter mit je 36 Karten: ein französisches Standardblatt von
AGMüller und das deutschschweizerische Blatt \enquote{Schaffhauser
Spielkarten Opti-Quatro} desselben Herstellers.

Wie gut das Modell leicht abweichende Kartendesigns oder Sondereditionen
erkennt, wurde nicht untersucht. Diese Grenze wird in
Abschnitt~\ref{sec:grenzen} erläutert.
